\begin{ouverture}[grandeurs]{13}{Les angles}{Partie 3 · Grandeurs et mesures}
{Comparer · Reporter · Sommes · L’angle droit vaut 90°}
\objectif{Comparer et reproduire des angles avec un calque.}
\objectif{Construire la somme ou un multiple d’angles donnés.}
\objectif{Utiliser la mesure de l’angle droit : 90 degrés.}
\prerequis{Angles comparés au CM1 et équerre.}
\end{ouverture}
\pagestyle{otMain}\essentiel\fontsize{11}{15}\selectfont
\phantomsection\label{cours-13-1}
\essSection{1}{Comparer et reproduire}
\begin{definition}Un angle est l’écartement entre deux côtés issus d’un même sommet. L’angle de sommet B dont les côtés passent par A et C se note $\widehat{ABC}$ : la lettre du sommet au milieu. La longueur dessinée des côtés ne change pas l’angle.\end{definition}
\label{fig-cours-13-1-cours-angle-calque}\figureGeo{cours-angle-calque}

\begin{savoirfaire}[Le calque, outil des angles]Pour comparer ou reproduire un angle : le décalquer, puis poser le calque sommet sur sommet, un côté sur un côté. L’angle le plus ouvert est le plus grand. Pour reporter, tracer le long des deux côtés du calque.\end{savoirfaire}
\begin{propriete}L’angle droit est le repère central : l’équerre le vérifie. Un angle peut être plus petit, égal ou plus grand qu’un angle droit ; au cours moyen, on travaille avec des angles saillants, au plus égaux à l’angle plat.\end{propriete}
\clearpage
\phantomsection\label{cours-13-2}
\essSection{2}{Sommes, multiples et degrés}
\begin{definition}Les angles se mesurent en degrés. Un angle droit mesure $90^\circ$. Deux angles droits mis côte à côte forment un angle plat : $90+90=180^\circ$. Le rapporteur, lui, sera utilisé au collège.\end{definition}
\begin{savoirfaire}[Juxtaposer pour additionner]Construire la somme de deux angles : reporter le premier, puis reporter le second contre un côté du premier, sommets superposés. Le double d’un angle s’obtient en le reportant deux fois.\end{savoirfaire}
\begin{exemple}La moitié d’un angle droit, obtenue par pliage, mesure $90:2=45^\circ$. Un angle droit et sa moitié mis bout à bout mesurent $90+45=135^\circ$. Les moitiés et les sommes donnent des mesures sans instrument.\end{exemple}
\begin{attention}Partager un angle en parts égales se fait par pliage du calque, pas à l’œil. Et une mesure annoncée se contrôle : un angle « de 90 degrés » doit accepter l’équerre exactement.\end{attention}

\clearpage\rubrique{Découverte}
\begin{blocLivre}{ex-13-1}\begin{exercice}[1]{Nommer et décrire}
Nommer l’angle marqué sur la figure, donner son sommet et ses côtés.
\label{fig-ex-13-1}\figureGeo{ang-1}
\end{exercice}\end{blocLivre}
\begin{blocLivre}{ex-13-2}\begin{exercice}[1]{Comparer au calque}
Ranger les trois angles dessinés du plus fermé au plus ouvert.
\label{fig-ex-13-2}\figureGeo{ang-2}
\end{exercice}\end{blocLivre}
\begin{blocLivre}{ex-13-3}\begin{exercice}[1]{Reproduire un angle}
Reproduire l’angle donné à l’aide d’un calque, à partir de la demi-droite déjà tracée.
\label{fig-ex-13-3}\figureGeo{ang-3}
\end{exercice}\end{blocLivre}
\begin{blocLivre}{ex-13-4}\begin{exercice}[1]{La chasse aux angles droits}
Avec l’équerre, repérer les angles droits du quadrilatère dessiné.
\label{fig-ex-13-4}\figureGeo{ang-4}
\end{exercice}\end{blocLivre}
\begin{blocLivre}{ex-13-5}\begin{exercice}[1]{90 degrés}
Un angle droit mesure 90 degrés. Que mesure un angle plat ? Et la moitié d’un angle droit ?
\end{exercice}\end{blocLivre}

\rubrique{Entraînement}
\begin{blocLivre}{ex-13-6}\begin{exercice}[2]{La somme de deux angles}
Construire un angle égal à la somme des deux angles donnés, puis dire s’il dépasse un angle droit.
\label{fig-ex-13-6}\figureGeo{ang-6}
\end{exercice}\end{blocLivre}
\begin{blocLivre}{ex-13-7}\begin{exercice}[2]{Le double}
Construire le double de l’angle donné. S’il mesure 35 degrés, que mesure son double ?
\label{fig-ex-13-7}\figureGeo{ang-7}
\end{exercice}\end{blocLivre}
\begin{blocLivre}{ex-13-8}\begin{exercice}[2]{La moitié par pliage}
Comment obtenir la moitié de l’angle dessiné sans instrument de mesure ? Si l’angle est droit, que mesure sa moitié ?
\label{fig-ex-13-8}\figureGeo{ang-8}
\end{exercice}\end{blocLivre}
\begin{blocLivre}{ex-13-9}\begin{exercice}[2]{Des mesures sans rapporteur}
Donner la mesure : d’un angle droit et demi ; de trois moitiés d’angle droit ; d’un angle plat moins un angle droit.
\end{exercice}\end{blocLivre}
\begin{blocLivre}{ex-13-10}\begin{exercice}[2]{Ranger des angles}
Ranger les quatre angles dessinés du plus petit au plus grand, en utilisant l’angle droit comme repère.
\label{fig-ex-13-10}\figureGeo{ang-10}
\end{exercice}\end{blocLivre}

\rubrique{Problèmes}
\begin{blocLivre}{ex-13-11}\begin{exercice}[3]{Les aiguilles de l’horloge}
À quelles heures piles les aiguilles d’une horloge forment-elles un angle droit ? Quel angle forment-elles à 6 h ?
\label{fig-ex-13-11}\figureGeo{ang-11}
\end{exercice}\end{blocLivre}
\begin{blocLivre}{ex-13-12}\begin{exercice}[3]{L’éventail}
Un éventail déplié forme un angle plat, partagé en 4 secteurs identiques. Quelle est la mesure de l’angle de chaque secteur ?
\label{fig-ex-13-12}\figureGeo{ang-13}
\end{exercice}\end{blocLivre}
\begin{blocLivre}{ex-13-13}\begin{exercice}[3]{Le toit de la cabane}
Les deux pans du toit forment ensemble un angle droit au faîte, et sont également inclinés. Quel angle chaque pan fait-il avec la ligne du faîte ?
\label{fig-ex-13-13}\figureGeo{ang-14}
\end{exercice}\end{blocLivre}
\begin{blocLivre}{ex-13-14}\begin{exercice}[3]{Des côtés trompeurs}
Zina affirme : « L’angle E est plus grand que l’angle F, ses côtés sont plus longs. » Corriger en comparant au calque.
\label{fig-ex-13-14}\figureGeo{ang-15}
\end{exercice}\end{blocLivre}
\begin{blocLivre}{ex-13-15}\begin{exercice}[3]{Le programme d’angles}
Exécuter : tracer une demi-droite ; reporter un angle droit ; contre lui, reporter la moitié d’un angle droit. Quelle est la mesure de l’angle total obtenu ?
\label{fig-ex-13-15}\figureGeo{ang-9}
\end{exercice}\end{blocLivre}

\clearpage
